
\paragraph{Full runs: MULTI-autoS.} The sweep curves are only a proxy (2000-NP cap), so we also
ran the rule for real. MULTI-autoS uses $\epsilon=c\,Q_{\nScaleBest}(S(\mathbf u_{\rm pilot}))$
with $c=\nScaleC$, the noise floor and the full budget; nothing else changes. On the main set
(where $c$ is in-sample) it is on average \nMainAutoSLoss\,dB below oracle MULTI, against
\nMainAutoLoss\,dB for MULTI-auto; it is within 0.5\,dB of the oracle on
\nMainAutoSWithinHalf/\nMainN{} problems (MULTI-auto: \nMainAutoWithinHalf), and its worst RE
loss drops from \nMainAutoRelLossMax\% to \nMainAutoSRelLossMax\%
(Table~\ref{tab:variants}, last row). The gain is concentrated at $\delta=2.5\%$, exactly where
$\kappa\hat\sigma^2$ overshoots.

For a held-out test we ran both automatic rules on the 12 Kodak images at the two noise levels
where they differ (no oracle needed). At $\delta=2.5\%$, MULTI-autoS is better on
\nKRHiWins/\nKRHiN{} problems, by \nKRHiGain\,dB ISNR on average (up to \nKRHiBestGain\,dB;
mean RE change \nKRHiRelRE\%). At $\delta=0.5\%$ the two rules are level
(\nKRLoGain\,dB on average, better on \nKRLoWins/\nKRLoN{}, worst \nKRLoWorstGain\,dB), and at
$\delta=1\%$, where $\kappa$ is effectively calibrated, they are also level
(\nISNRKodakAutoS{} vs.\ \nISNRKodakAuto\,dB, Table~\ref{tab:kodak}). So the scale-aware rule
removes the high-noise failure of $\kappa\hat\sigma^2$ with no systematic loss elsewhere, at the cost of
one pilot solve (about 40 FFTs). We keep $\kappa\hat\sigma^2$ as ``MULTI-auto'' in the main
tables because it was fixed before these experiments; MULTI-autoS is reported alongside it.
